\section{Structural Framework}
  \label{sec:structural-framework}

  \subsection{Projection Fiber and Effective $U(1)$ Structure}
    \label{subsec:projection-fiber}

    The present mechanism is formulated within a minimal structural setting.
    We assume that the effective electromagnetic gauge symmetry arises as
    the $U(1)$ fiber symmetry of a projective structure
    \begin{equation}
      \Pi \cong S^3 ,
    \end{equation}
    whose Hopf fibration defines a base manifold and a circular fiber.
    The fiber coordinate
    \begin{equation}
      \theta \in [0,2\pi)
    \end{equation}
    plays the role of an effective phase variable.

    The electromagnetic potential $A$ is identified with the connection
    associated with this $U(1)$ fiber in the projected description.
    It is not an additional gauge field introduced independently,
    but the connection of the fiber as it appears in the effective theory.
    Gauge covariance therefore fixes the only admissible covariant gradient
    for the fiber phase,
    \begin{equation}
      D\theta \equiv \nabla \theta - \frac{q}{\hbar} A ,
      \label{eq:covariant-theta}
    \end{equation}
    where $q$ denotes the effective topological charge
    of the configuration under consideration.

    No further gauge postulate is introduced.
    The $U(1)$ structure is entirely inherited from the fiber geometry.

  \subsection{Topological Excitations and Composite Classes}
    \label{subsec:topological-excitations}

    Charged excitations correspond to topologically non-trivial
    configurations of the fiber phase.
    A single elementary excitation carries winding number
    \begin{equation}
      w = 1 ,
    \end{equation}
    and couples to the connection with effective charge $q=e$.

    A bound composite of two such excitations forms a distinct
    topological class characterized by
    \begin{equation}
      w = 2 .
    \end{equation}
    This class is not simply the superposition of two independent windings.
    It corresponds to a globally coherent configuration in which
    the fiber phase interpolates continuously between the two cores,
    subject to a raccordement constraint.
    Separating the cores introduces a phase defect
    whose energetic cost defines a stability threshold.

    For the $w=2$ class, the covariant gradient becomes
    \begin{equation}
      D\theta = \nabla \theta - \frac{2e}{\hbar} A .
      \label{eq:covariant-theta-2e}
    \end{equation}
    This composite degree of freedom therefore carries effective charge $2e$
    and supports a collective phase variable $\theta$.

  \subsection{Projective Gap and Stability Criterion}
    \label{subsec:projective-gap}

    We define the \emph{projective gap} $\Delta_\Pi$
    as the minimal excitation energy required
    to destabilize the $w=2$ composite class.
    Operationally, this corresponds to the lowest positive eigenvalue
    of a linearized relaxation operator
    restricted to the $U(1)$ fiber sector
    around the composite configuration,
    \begin{equation}
      \Delta_\Pi \equiv
      \lambda_{\min}
      \left(
        \mathcal{L}_\Pi \big|_{w=2}
      \right) .
      \label{eq:projective-gap}
    \end{equation}

    Physically, $\Delta_\Pi$ is the minimal energy cost
    to either separate the composite into two $w=1$ excitations
    or introduce a local phase slip that alters
    the global winding class.
    The existence of a finite $\Delta_\Pi$
    is the structural condition for composite stability.

    Two distinct mechanisms may generate such stability:

    \begin{itemize}
      \item In conventional low-$T_c$ systems,
      a bosonic mediator lowers the effective raccordement cost,
      stabilizing the composite through an attractive channel.
      \item In strongly correlated systems,
      stability arises from the reduction of real-space frustration
      in the admissibility constraints of the fiber,
      without invoking a separate mediator.
    \end{itemize}

    The remainder of this work shows that
    once a stable $w=2$ composite exists,
    global phase locking of the fiber
    naturally generates the London and Ginzburg--Landau structures,
    flux quantization, and Meissner screening.

  \subsection{Free-Energy Functional from Phase Coherence}
    \label{subsec:free-energy-functional}

    When a macroscopic density of $w=2$ composites
    occupies the same coherent phase sector,
    we define the order parameter
    \begin{equation}
      \Psi = |\Psi| e^{i\theta} ,
    \end{equation}
    where $|\Psi|^2$ measures the density of composites
    participating in coherent projection.

    The leading contribution to the free energy
    must be quadratic in the only gauge-covariant derivative available,
    \begin{equation}
      \mathcal{F}
      \supset
      \frac{\hbar^2 \rho_s}{2 m^*}
      \left(
        \nabla \theta - \frac{2e}{\hbar} A
      \right)^2 ,
      \label{eq:london-structure}
    \end{equation}
    where $\rho_s$ is the phase stiffness
    and $m^*$ the effective inertia of the composite mode.

    Equation~\ref{eq:london-structure}
    is structurally identical to the London
    and Ginzburg--Landau gradient terms.
    It arises here from the energetic cost
    of breaking global fiber coherence,
    not from an assumed symmetry-breaking potential.

    The Meissner effect, flux quantization,
    and the effective mass of the gauge field
    follow directly from Eq.~\ref{eq:london-structure}.
    These consequences will be developed explicitly
    in the following sections.
